\documentclass[11pt,a4paper]{article}

\usepackage[margin=1in]{geometry}
\usepackage[T1]{fontenc}
\usepackage[utf8]{inputenc}
\usepackage{lmodern}
\usepackage{microtype}
\usepackage{amsmath,amssymb,mathtools}
\usepackage{booktabs}
\usepackage{graphicx}
\usepackage{float}
\usepackage[hidelinks]{hyperref}

\graphicspath{{../demos/report_outputs/exercise_1/}{Balint/demos/report_outputs/exercise_1/}}

\title{Final Report}
\author{}
\date{}

\begin{document}

    \maketitle


    \section*{Exercise 1: Entanglement Definitions and Probes}

    This section studies entanglement diagnostics in the open one-dimensional transverse-field Ising model using the project-owned exact-diagonalization workflow. The module structure is organized to reflect the same high-level API used in NetKet: graph, Hilbert space, operator, exact solver, sampler, and variational-state logic are kept conceptually separate so that later Monte Carlo workflows can reuse the same building blocks and so that a specific NQS workflow is simple to set up. For the one-dimensional TFIM, the expected asymptotic behavior is area-law saturation in the gapped phases and logarithmic block-entropy growth at criticality. The chosen field values are as follows: $h=0.5$ is the ordered reference point, $h=1.5$ is the disordered reference point, and $h=0.95$, $h=1.0$, and $h=1.05$ form the near-critical cluster used to describe the crossover region around the critical point.

    \subsection*{1/a Exact Diagonalization Setup}

    The exact-diagonalization path starts from the Hilbert-space definition provided by \texttt{SpinHilbert}. The current implementation is restricted to spin-$1/2$ systems, with configurations encoded as binary bitstrings with entries in $\{0,1\}$. The Hilbert-space object provides the total dimension $2^N$ and implements the map between basis states and integer indices. Site $0$ is stored as the least-significant bit, so the basis index of a configuration is
    \begin{equation}
        n = \sum_i s_i 2^i .
    \end{equation}
    For $L=16$, the Hilbert-space dimension is $2^{16}=65{,}536$, which is still small enough for exact diagonalization and explicit reshaping of the ground-state vector.

    The lattice geometry is supplied by the open \texttt{Chain1D} graph object, which provides a deterministic site ordering together with the nearest-neighbor bond list used in the Hamiltonian. This keeps the basis convention separate from the interaction pattern and makes the same infrastructure reusable for both one-dimensional chains and two-dimensional square lattices elsewhere in the project.

    The Hamiltonian is assembled through the TFIM operator builder, which constructs the operator as a sum of local $\sigma_i^z \sigma_j^z$ interaction terms over graph edges and on-site $\sigma_i^x$ field terms,
    \begin{equation}
        H = -J \sum_{\langle ij \rangle} \sigma_i^z \sigma_j^z - h \sum_i \sigma_i^x .
    \end{equation}
    The operator representation stores these contributions as local terms and exposes the connected matrix elements $H_{\sigma',\sigma}$ generated by acting on a basis configuration $\lvert \sigma \rangle$. This is important not only for exact diagonalization but also for the variational Monte Carlo estimator used later, where the local energy is written as
    \begin{equation}
        E_{\mathrm{loc}}(\sigma) = \sum_{\sigma'} H_{\sigma,\sigma'} \frac{\psi(\sigma')}{\psi(\sigma)} ,
    \end{equation}
    and the variational energy is obtained by averaging $E_{\mathrm{loc}}(\sigma)$ over configurations sampled from $\lvert \psi(\sigma) \rvert^2$.

    The exact solver then iterates over the operator matrix elements, builds a sparse COO matrix in coordinate-wise form, converts it to CSR format, and calls \path{scipy.sparse.linalg.eigsh} to obtain the lowest eigenpair using a Lanczos-based sparse eigensolver. The global phase of the ground state is fixed afterwards so that the amplitude with the largest magnitude is real. This sparse construction is preferable to forming a dense $65536 \times 65536$ matrix and already matches the connected-elements view that is reused later in the sampled VMC calculations.

    As a concrete example, the discussion below uses the half-chain split $A|B$ with $\lvert A \rvert = \lvert B \rvert = 8$. The Hilbert space then factorizes as
    \begin{equation}
        \mathcal{H} = \mathcal{H}_A \otimes \mathcal{H}_B ,
    \end{equation}
    with dimensions $2^{\lvert A \rvert}$ and $2^{\lvert B \rvert}$. Numerically, this means reshaping the exact state vector into a $2^{\lvert A \rvert} \times 2^{\lvert B \rvert}$ amplitude matrix $\psi(s_A,s_B)$, which is the key step for calculating reduced density matrices and Schmidt decompositions.

    \subsection*{1/b Von Neumann Entropy and the Reduced Density Matrix}

    For an exact pure state $\lvert \psi \rangle$, the reduced density matrix of subsystem $A$ is
    \begin{equation}
        \rho_A = \mathrm{tr}_B \lvert \psi \rangle \langle \psi \rvert ,
    \end{equation}
    and the von Neumann entropy is
    \begin{equation}
        S_1(A) = -\mathrm{tr}(\rho_A \log \rho_A) .
    \end{equation}
    The eigenvalues of $\rho_A$ are the squared Schmidt coefficients across the cut. When considering matrix product states, they quantify how much weight is carried by each Schmidt sector and therefore determine how compressible the state is across that bond. A rapidly decaying spectrum indicates that a small bond dimension would retain most of the state weight, while a broader spectrum implies weaker compressibility.

    The exact half-chain entropies for the retained field values are listed in Table~\ref{tab:exact-entropies}.

    \begin{table}[H]
        \centering
        \begin{tabular}{llrr}
            \toprule
            $h$  & Role                                 & $S_1(A)$ & $S_2(A)$ \\
            \midrule
            0.5  & ordered reference                    & 0.6990   & 0.6944   \\
            0.95 & near-critical point, ordered side    & 0.7762   & 0.6458   \\
            1.0  & critical point                       & 0.7501   & 0.5734   \\
            1.05 & near-critical point, disordered side & 0.6970   & 0.4842   \\
            1.5  & disordered reference                 & 0.3074   & 0.1419   \\
            \bottomrule
        \end{tabular}
        \caption{Exact half-chain entropies for the five retained TFIM field values on the open $L=16$ chain.}
        \label{tab:exact-entropies}
    \end{table}

    For the one-dimensional TFIM, the expected asymptotic distinction is that gapped phases obey area-law saturation while the critical point shows logarithmic growth with subsystem size. The present $L=16$ data do not cleanly expose that asymptotic separation: the half-chain entropy is largest at the near-critical ordered-side point $h=0.95$, while $h=1.0$ and $h=1.05$ remain numerically close to it. The useful interpretation is therefore that this finite open chain mainly resolves a crossover region around criticality rather than a clean thermodynamic-limit scaling law.

    The entanglement spectrum shown in Figure~\ref{fig:spectrum} supports the same interpretation.

    \begin{figure}[H]
        \centering
        \includegraphics[width=0.82\linewidth]{exercise_1_entanglement_spectrum.png}
        \caption{Leading Schmidt weights across the half-chain cut for the five retained field values.}
        \label{fig:spectrum}
    \end{figure}

    The spectrum becomes more concentrated as the field is increased toward the disordered regime. The Schmidt-compressibility data make this quantitative: keeping only the leading Schmidt value retains about $50.05\%$ of the weight at $h=0.5$, $64.74\%$ at $h=0.95$, $70.45\%$ at $h=1.0$, $75.91\%$ at $h=1.05$, and $93.02\%$ at $h=1.5$. The retained weight therefore increases away from the near-critical region, so truncation is easiest deep in the disordered phase and hardest on the ordered side and near criticality.

    \subsection*{1/c Renyi-2 Entropy and the SWAP Identity}
    $S_1$ remains primarily an exact-benchmark observable in the present context. Evaluating the von Neumann entropy requires explicit access to $\rho_A$ or all of its eigenvalues, and that becomes exponentially expensive with subsystem size. For small exact-diagonalization benchmarks it is straightforward; for a sampled neural quantum state it is not generally an efficient default observable, as it would reintroduce the same exponential wall that the variational ansatz is meant to avoid.

    The Renyi entropies are defined by
    \begin{equation}
        S_\alpha(A) = \frac{1}{1-\alpha} \log \mathrm{tr}(\rho_A^\alpha) .
    \end{equation}
    For $\alpha=2$, this reduces to
    \begin{equation}
        S_2(A) = -\log \mathrm{tr}(\rho_A^2) ,
    \end{equation}
    so the relevant quantity is the purity $\mathrm{tr}(\rho_A^2)$. The extremal limits are simple and useful checks: for a globally pure state that factorizes across the $A|B$ cut, $\rho_A$ is pure, so $\mathrm{tr}(\rho_A^2)=1$ and $S_2(A)=0$; for a maximally mixed reduced density matrix on a subsystem of dimension $d_A$, one has $\mathrm{tr}(\rho_A^2)=1/d_A$ and $S_2(A)=\log d_A$.

    The exact identity used later in the neural-quantum-state setting is
    \begin{equation}
        \mathrm{tr}(\rho_A^2)
        =
        \langle \psi \otimes \psi \rvert
        \mathrm{SWAP}_A
        \lvert \psi \otimes \psi \rangle .
    \end{equation}
    Writing the amplitudes as $\psi(s_A,s_B)$ gives
    \begin{equation}
        \rho_A(s_A,s_A') = \sum_{s_B} \psi(s_A,s_B)\psi^*(s_A',s_B) ,
    \end{equation}
    so that
    \begin{equation}
        \mathrm{tr}(\rho_A^2)
        =
        \sum_{s_A,s_A',s_B,s_B'}
        \psi(s_A,s_B)\psi^*(s_A',s_B)\psi(s_A',s_B')\psi^*(s_A,s_B') .
    \end{equation}
    This is exactly the replica expectation value of the operator that swaps the $A$ degrees of freedom between the two copies. Intuitively, the SWAP test asks how compatible two independent replicas remain after their subsystem-$A$ configurations are exchanged. If the state factorizes across the $A|B$ cut, swapping $A$ leaves the two-copy amplitude unchanged and the expectation is exactly $1$, so the subsystem is pure and the entropy ($log(1)$) is 0. If $A$ is strongly entangled with $B$, the swapped configurations are less compatible with the original amplitudes, so the expectation decreases, reflecting a more mixed reduced state on $A$ and increased entropy. The same construction can be generalized to $\alpha>2$ by introducing more replicas and a cyclic permutation. This is technically possible, but in the current implementation the estimator is already noisy enough that higher-$\alpha$ probes would not be practically useful even if an efficient estimator were written down.

    The SWAP-check data confirm the exact identity numerically for every retained field value. At $h=1.0$, for example, the purity from $\mathrm{tr}(\rho_A^2)$ is $0.5636148913$, while the purity reconstructed from $\exp(-S_2)$ is $0.5636148913$ to machine precision. The same agreement holds at $h=0.5$, $0.95$, $1.05$, and $1.5$, which validates the Renyi-2 route used later in the sampled neural-quantum-state setting.

    Because $S_2$ depends only on a low-order trace of $\rho_A$, it is much more practical for Monte Carlo estimation than $S_1$. One still expects the estimator to become harder for larger subsystems, but it avoids the need to reconstruct the full reduced density matrix or all of its eigenvalues.

    \subsection*{1/d Interpreting $S_1$ and $S_2$}

    For gapped one-dimensional ground states, entanglement is expected to obey an area law: for a fixed boundary cut, $S_1(A)$ and $S_2(A)$ should remain $\mathcal{O}(1)$ instead of growing linearly with subsystem size. A volume-law state would instead produce entropy proportional to $\lvert A \rvert$. Long-range entanglement is not simply ``large entropy''; it refers to nontrivial global entanglement structure that cannot be reduced to a short-range boundary contribution alone.

    The subsystem-size scan is shown in Figure~\ref{fig:scaling}.

    \begin{figure}[H]
        \centering
        \includegraphics[width=\linewidth]{exercise_1_entropy_scaling.png}
        \caption{Von Neumann and Renyi-2 entropy versus subsystem size for the five retained field values.}
        \label{fig:scaling}
    \end{figure}

    The main finite-size message from this plot is not a clean asymptotic scaling law, but a crossover comparison. The disordered reference point $h=1.5$ stays low across subsystem sizes and is the clearest area-law-like anchor in the retained data, which is what is expected for a gapped non-critical phase in the one-dimensional TFIM. At the same time, the finite system does not cleanly separate the close-to-critical points from the critical point itself: the near-critical trio $h=0.95$, $1.0$, and $1.05$ all show enhanced growth relative to $h=1.5$, with the largest finite-size entropies appearing slightly on the ordered side of criticality rather than exactly at $h=1.0$. The plot therefore serves mainly as a finite-size crossover window rather than as a clean asymptotic scaling test.

    The heuristic log-fit summary should be read cautiously. Its $R^2$ values are high for $h=0.95$, $1.0$, and $1.05$, which is consistent with enhanced finite-size growth around the crossover region, but the system is too small for strong claims about asymptotic critical scaling. The log-fit summary is therefore best read only as a compact finite-size diagnostic.

    Truncating the smallest Schmidt coefficients affects $S_1$ more strongly than $S_2$ because the two probes weight the entanglement spectrum differently. The von Neumann entropy contains the factor $-\lambda \log \lambda$, so even small Schmidt weights contribute. By contrast, $S_2$ depends on $\sum_i \lambda_i^2$, which emphasizes the dominant Schmidt sectors and suppresses the tail. This is why $S_1$ is the more tail-sensitive diagnostic, while $S_2$ is the more practical probe when moving to sampled neural quantum states.

    \subsection*{Conclusion}

    Exact diagonalization on the open $L=16$ TFIM chain provides a compact reference point for the entanglement probes used later in the project. The computational-basis convention, explicit graph structure, local-term Hamiltonian construction, and sparse Lanczos solve together define a consistent benchmark path that is also aligned with the connected-matrix-elements view required later in VMC.

    The reduced-density-matrix eigenvalues define the Schmidt spectrum across the half-chain cut and therefore quantify compressibility directly in matrix-product-state language. Across the retained field values, truncation becomes easier away from the critical regime, with the disordered point $h=1.5$ the most compressible case and the near-critical points remaining less compressible.

    The SWAP check verifies the Renyi-2 identity at machine precision for all five field values. This matters because $S_2$ remains a practical entanglement probe in the neural-quantum-state setting, whereas $S_1$ generally requires explicit reduced-density-matrix reconstruction. For the next exercises, this motivates carrying Renyi entropies into the neural-quantum-state setting.

    The finite-size scan is most naturally interpreted by taking $h=1.5$ as the low-entanglement area-law reference and $h=0.95$, $1.0$, and $1.05$ as the near-critical cluster for discussing crossover behavior. On this system size, however, the close-to-critical points and the critical point do not separate cleanly, and the largest half-chain entropies appear at $h=0.95$ rather than exactly at $h=1.0$. The contrast between $S_1$ and $S_2$ is nevertheless clear: $S_1$ is more sensitive to the tail of the entanglement spectrum, while $S_2$ is more robustly tied to the dominant Schmidt sectors.

    It's also worth keeping in mind that clearer scaling laws should emerge only at larger system sizes.

\end{document}
